\scanpage{70}
Supposons \unsure{en|de} plus

d) toute flèche inversible de $\uline{\Sigma}$ est une identité
(p.ex. $\uline{\Sigma}$ est une catégorie définie par un ens
\uline{ordonné}).

Alors on conclut aisément que l'ens. pré-ordonné $I$ est
\uline{ordonné}, et que tte flèche de $\uline{\Sigma}$ qui n'est pas
une identité i.e. $u\in\Sigma^{*}(1)$, son image dans $I(1)$
n'est pas une identité;
\deleted{\unsure{i.e.}} Si $\uline{\Sigma}$ est associé à un
\deleted{\unsure{ens. ord.}}\inserted{espace top.}
ordonné, $\Sigma=\Sigma_{(0)}$, cela signifie que
$\uline{\Sigma}\longrightarrow\uline{I}$ est \inserted{une appl.}
\uline{strictement} croissante.

Revenons alors à un \deleted{catég} ens. ordonné topologique
\[
\uline{\Sigma}=(\Sigma=\Sigma_{(0)},\quad
\underset{\underset{\Sigma\times\Sigma}{\cap}}{\Sigma_{(1)}}
=\delta(\Sigma_{(0)})\amalg\Sigma^{<}(1))
\]
satisfaisant aux conditions a) b) c) et 1°) de la page (27)
(correspondant \inserted{\unsure{moralement}} à un déploiement d'une
stratification locale équisingulier), et tel que les comp. connexes
de $\Sigma$ \deleted{sont} de dim finie.
Considérons alors \unsure{la|le} \deleted{\unsure{factorisation}}\inserted{morphisme}
canonique $\uline{\Sigma}\longrightarrow\uline{I}$ avec $I=\pi_0(\Sigma)$,
et supposons de plus que pour $i<j$ ($i,j\in I$)
\[
\Sigma_{ij}\longrightarrow\Sigma_j\quad\hbox{soit \uline{mono} i.e. plongement fermé.}
\]
On est alors sous les conditions de p. 28,
données a) b) c)
$(\Sigma_i,\Sigma_{ij},\ s_{ij},b_{ij}:$
\diagram{\begin{tikzcd}[column sep=large,row sep=small]
&\Sigma_i\\
\Sigma_{ij}\arrow[ur]\arrow[dr]&\\
&\Sigma_j
\end{tikzcd}}$)$,
\mspage{69}
